\documentclass[12pt]{article}
\usepackage[utf8]{inputenc}
\usepackage[T1]{fontenc}
\usepackage[margin=1in]{geometry}
\usepackage{setspace}
\usepackage{xcolor}
\usepackage[colorlinks=true, urlcolor=blue!60!black, linkcolor=blue!60!black]{hyperref}
\setlength{\parindent}{2em}
\setlength{\parskip}{1.25em}
\renewcommand{\baselinestretch}{1.0}
\begin{document}
\begin{center}
{\LARGE \textit{Cedar Point Nursery v.\ Hassid}, 594 U.S.\ 139 (2021)\par}
\vspace{0.8em}
{\large Case Brief --- Four-Page Edition\par}
\end{center}
\vspace{1em}

\section{Memory Jogger}

A California rule letting union organizers enter farms for up to three hours a day, 120 days a year, was held 6--3 to be a \emph{per se} physical taking: a regulation that appropriates the owner's right to exclude for third parties requires just compensation, however intermittent the access, subject to limited exceptions (\href{https://supreme.justia.com/cases/federal/us/594/20-107/}{Justia}).

\section{Detailed Case Facts}

Under California's Agricultural Labor Relations Act of 1975, the Agricultural Labor Relations Board adopted a regulation giving labor organizations a ``right to take access'' to an agricultural employer's premises to meet employees and solicit support. After written notice, organizers could enter for up to three hours a day---an hour before work, at lunch, and after work---during four 30-day periods a year; interference was an unfair labor practice. Two growers whose workers lived off-site challenged the rule. Cedar Point Nursery, a strawberry grower with over 400 seasonal and about 100 full-time workers, saw United Farm Workers organizers enter its trim shed at 5:00 a.m.\ in October 2015 without notice, using bullhorns. Fowler Packing, a Fresno grape and citrus grower, blocked organizers in July 2015, prompting a charge the union later withdrew.

\section{Procedural History}

The growers sued Board officials, claiming the regulation appropriated an uncompensated easement. The district court dismissed, holding that because access was neither permanent nor continuous, \textit{Penn Central} balancing applied, which the growers had not pleaded. A divided Ninth Circuit affirmed, limiting \emph{per se} treatment to access ``24 hours a day, 365 days a year.'' Rehearing en banc was denied over Judge Ikuta's dissent, joined by seven judges. The Supreme Court reversed on June 23, 2021.

\section{Judicial Votes}

\textbf{6--3}, reversed and remanded.
\begin{itemize}
\item \textbf{Majority:} Chief Justice Roberts, joined by Justices Thomas, Alito, Gorsuch, Kavanaugh, and Barrett.
\item \textbf{Concurring:} Justice Kavanaugh.
\item \textbf{Dissenting:} Justice Breyer, joined by Justices Sotomayor and Kagan.
\end{itemize}

\section{Holding}

The access regulation appropriates a right to physically invade the growers' property, transferring to third parties the owners' right to exclude, and is therefore a \emph{per se} physical taking requiring just compensation, whether the access is intermittent or temporary. Isolated trespasses, entries consistent with background property law, and reasonable inspection conditions on government benefits are not takings.

\section{Analysis of Opinions}

\textbf{Majority (Roberts): appropriation versus use restriction.} Restrictions on how owners use their property get \textit{Penn Central} balancing; physical appropriations for the government or a third party---by eminent domain, seizure (\textit{Pewee Coal}), or regulation (\textit{Horne})---require payment, and \textit{Penn Central} has no place.

\textbf{Majority (Roberts): the right to exclude.} The right to exclude is among the most treasured strands in the bundle (\textit{Loretto}; \textit{Kaiser Aetna}), and authorized invasions ``by plane, boat, cable, or beachcomber'' are takings (\textit{Causby}; \textit{Nollan}). Duration bears on the amount of compensation, not liability (\textit{United States v.\ Dow}), and the regulation's mismatch with a state-law easement is irrelevant. \textit{PruneYard} involved a shopping center open to the public; \textit{Babcock \& Wilcox} was a statutory NLRA case.

\textbf{Majority (Roberts): three limits.} (1) Isolated invasions not made under a granted right of access remain torts. (2) Entries consistent with background privileges---nuisance abatement (\textit{Lucas}), necessity, reasonable searches---take nothing the owner had. (3) Government may condition permits and licenses on reasonable health and safety inspections that satisfy \textit{Nollan} and \textit{Dolan}.

\textbf{Kavanaugh (concurring).} \textit{Babcock \& Wilcox} independently recognized employers' right to exclude organizers absent necessity, such as workers living on company property.

\textbf{Breyer, joined by Sotomayor and Kagan (dissenting).} The rule regulates rather than appropriates the right to exclude and creates no estate or easement. Only permanent occupations are \emph{per se} takings; temporary access, about 4\% of the year, belongs under \textit{Penn Central}. The new rule threatens many inspections, and under \textit{Knick} California should be able to avoid an injunction by paying compensation.

\section{Examples: Future Applications}

\begin{enumerate}
\item \textbf{Same-Side: Wildlife Surveys.} A statute letting nonprofits enter fenced private timberland four hours a week for surveys appropriates a right of entry; duration affects only compensation.
\item \textbf{Same-Side: Gated-Community Canvassers.} Forcing private gated communities to admit political canvassers grants third parties entry; \textit{PruneYard} is confined to property open to the public.
\item \textbf{Opposite-Side: Licensing Inspections and Zoning.} Food-safety inspections as a license condition fit the third exception, and a 500-foot zoning setback restricts use, triggering \textit{Penn Central}.
\item \textbf{Opposite-Side: One-Time Intrusion.} A highway crew that mistakenly stores gravel on a field for two days commits an isolated trespass, not a \emph{per se} taking.
\item \textbf{Fence-Sitter: Supermarket Organizers.} A statute giving organizers access to a supermarket's private sidewalk tests how far \textit{PruneYard} survives for property open to shoppers but not solicitors.
\end{enumerate}

\section{Critique}

\textbf{Remedial paradox.} As Justice Breyer noted, the Takings Clause requires compensation rather than forbidding takings, yet the growers sought an injunction, leaving \textit{Cedar Point}'s interaction with \textit{Knick} unresolved. \textbf{Fuzzy exceptions.} Lines between isolated torts and recurring takings, which common-law privileges count, and what counts as a government benefit will generate litigation over environmental, labor, and housing rules. \textbf{Fragile \textit{PruneYard}.} Limiting it to businesses generally open to the public invites challenges by retail venues seeking to exclude non-customers.

\section{Key Quotations}

\begin{quote}
``The essential question is not, as the Ninth Circuit seemed to think, whether the government action at issue comes garbed as a regulation (or statute, or ordinance, or miscellaneous decree). It is whether the government has physically taken property for itself or someone else---by whatever means---or has instead restricted a property owner's ability to use his own property.''
\end{quote}
\noindent --- Chief Justice Roberts, for the Court.

\begin{quote}
``The access regulation appropriates a right to invade the growers' property and therefore constitutes a per se physical taking. \ldots{} There is no reason the law should analyze an abrogation of the right to exclude in one manner if it extends for 365 days, but in an entirely different manner if it lasts for 364.''
\end{quote}
\noindent --- Chief Justice Roberts, for the Court.

\begin{quote}
``We cannot agree that the right to exclude is an empty formality, subject to modification at the government's pleasure. On the contrary, it is a `fundamental element of the property right,' \textit{Kaiser Aetna}, 444 U.S., at 179--180, that cannot be balanced away.''
\end{quote}
\noindent --- Chief Justice Roberts, for the Court.

\end{document}
